% ─────────────────────────────────────────────────
%  CHAPITRE 4 – PRESENTATION DES RESULTATS ET ANALYSE CRITIQUE
% ─────────────────────────────────────────────────
\chapter{Présentation des résultats et analyse critique}
\markboth{Présentation des résultats et analyse critique}{Présentation des résultats et analyse critique}

\section*{Introduction}
\phantomsection
\addcontentsline{toc}{section}{Introduction}


Après l'analyse des besoins et la conception de la solution, ce chapitre présente la réalisation de la plateforme I-SIB. Il s'intéresse aux fonctionnalités mises en œuvre dans les modules Relation Client et SGI, en s'appuyant sur plusieurs scénarios représentatifs du périmètre étudié.

Les exemples présentés ont été réalisés dans l'environnement de validation. Les données utilisées pour la partie Relation Client sont des données de démonstration, tandis que les informations boursières reposent sur les cotations BRVM \cite{brvm_data_public}, issues d'une source publique mise à jour quotidiennement.

Le chapitre se termine par un retour sur les principaux apports de la solution ainsi que sur les améliorations qui pourront être envisagées par la suite.


% ─────────────────────────────────────────────────
\section{Résultats fonctionnels obtenus}

\subsection{Centralisation de la connaissance client}

La réalisation a permis de centraliser les informations utiles au suivi
commercial. Le chargé de clientèle dispose d'un espace unique pour consulter son
portefeuille, suivre les statuts des clients, repérer les niveaux de risque et
accéder rapidement aux informations nécessaires à la relation client.

La vue de portefeuille retenue pour illustrer ce résultat est présentée à la
figure~\ref{fig:portefeuille_clients}. Elle montre l'organisation générale de
l'écran autour des indicateurs de synthèse, des outils de recherche et de la
liste des clients suivis.

\begin{figure}[H]
    \centering
    \MemoireFigureLarge{images/contenu/chapitre-4/portefeuille_clients_annote.png}
    \caption{Vue de portefeuille client utilisée pour le suivi commercial}
    \label{fig:portefeuille_clients}
\end{figure}

\noindent
Le repère \textbf{1} indique les principaux indicateurs de synthèse du
portefeuille, notamment le nombre de clients, les clients actifs, les nouvelles
acquisitions et le taux de rétention. Le repère \textbf{2} correspond à la zone
de recherche et aux filtres utilisés pour retrouver rapidement un client ou un
segment précis. Le repère \textbf{3} présente la liste détaillée des clients
suivis, avec les informations utiles au suivi commercial : statut, niveau de
risque, comptes, produits, solde et indicateur de risque de départ.

Cette vue regroupe les éléments nécessaires au suivi quotidien du portefeuille
et à l'identification des dossiers nécessitant une attention particulière.

\subsection{Suivi commercial et gestion des prospects}

La plateforme prend également en charge la phase amont de la relation client à
travers la gestion des prospects. Le prospect peut être enregistré, suivi dans
un pipeline commercial et associé à un score initial permettant de prioriser les
actions du chargé de clientèle. Ce résultat s'inscrit dans le prolongement du
scénario de création d'un prospect présenté dans la conception détaillée.

La gestion des prospects s'inscrit dans la continuité du suivi client. Les
informations saisies peuvent alimenter le dossier commercial et préparer la
conversion éventuelle du prospect en client. Dans cette version, le score
constitue un indicateur d'aide au suivi ; son évaluation devra être consolidée
avec des données réelles.

\subsection{Consultation des instruments et aide à la décision SGI}

Le second ensemble de résultats concerne le module SGI, c'est-à-dire les
fonctionnalités liées aux opérations d'investissement. La plateforme permet de
consulter des instruments référencés sur la BRVM \cite{brvm_actions}, notamment les actions et les
obligations. L'utilisateur peut accéder aux principaux
éléments nécessaires à la prise de décision : symbole, dénomination, prix,
variation, volume et classement des hausses ou baisses.

Dans l'environnement de validation, les cotations proviennent d'une source
publique tierce \cite{brvm_data_public} qui publie chaque jour les données de la
BRVM \cite{brvm_actions}, en attendant un accès direct au flux de marché
officiel. Cette phase permet de vérifier le parcours de consultation, le
filtrage et la lecture des instruments avant un branchement ultérieur sur ce
flux.

L'écran de cotation présenté à la figure~\ref{fig:cotations_brvm} montre
l'organisation de cette consultation autour des classements de titres, des
filtres de navigation et du tableau des instruments disponibles.

\begin{figure}[H]
    \centering
    \MemoireFigureLarge{images/contenu/chapitre-4/cotations_brvm_annote.png}
    \caption{Consultation d'instruments BRVM (source publique tierce, données quotidiennes)}
    \label{fig:cotations_brvm}
\end{figure}

\noindent
Le repère \textbf{1} indique le classement des plus fortes hausses, c'est-à-dire
les titres dont la variation est la plus favorable sur la période observée. Le
repère \textbf{2} signale les plus fortes baisses sur la même période. Le repère
\textbf{3} correspond aux filtres et aux onglets permettant de naviguer entre
indices, actions et obligations. Le repère \textbf{4} présente enfin le tableau
des instruments, avec les informations de marché utiles avant une décision
d'achat ou de vente.


% ─────────────────────────────────────────────────
\section{Analyse critique de la réalisation}

Cette section met en perspective les résultats obtenus au regard des attentes
fonctionnelles, de la qualité logicielle \cite{iso25010} et des exigences de
sécurité applicative \cite{owasp_asvs}.

\subsection{Évaluation de la dégradation gracieuse du scoring}

Nous avons mené une expérimentation pour tester la dégradation gracieuse du
scoring. Le principe consiste à créer 40 prospects via la passerelle
(\texttt{POST /api/crm/prospects}), une première fois avec le service de scoring
actif, puis une seconde fois après son arrêt. Le tableau~\ref{tab:eval_degradation}
présente les résultats obtenus. Pour chaque requête, le code HTTP et le temps de
réponse sont relevés, les appels étant espacés afin de respecter la limitation
de débit de la passerelle.

\begin{table}[h]
    \centering
    \renewcommand{\arraystretch}{1.4}
    \setlength{\tabcolsep}{10pt}
    \begin{tabular}{|l|c|c|c|c|}
        \hline
        \textbf{Condition} & \textbf{Taux de succès} & \textbf{Moyenne} & \textbf{p50} & \textbf{p95} \\
        \hline
        Scoring actif   & 40/40 (100\,\%) & 82\,ms   & 47\,ms   & 67\,ms \\
        \hline
        Scoring arrêté  & 40/40 (100\,\%) & 1593\,ms & 1544\,ms & 2034\,ms \\
        \hline
    \end{tabular}
    \caption{Création de prospect avec et sans service de scoring
        (40 requêtes par condition).}
    \label{tab:eval_degradation}
\end{table}

Dans les deux conditions, les 40 créations aboutissent~: la dégradation gracieuse
est totale.
Lorsque le service de scoring est indisponible, la création d'un prospect reste possible. Dans ce cas, le système attribue un score par défaut et la requête se termine normalement avec un code de réponse 201. Le fonctionnement du module de création ne dépend donc pas directement de la disponibilité du service de scoring.

La principale différence observée concerne le temps de réponse. Avec le service de scoring actif, la création d'un prospect s'effectue en quelques dizaines de millisecondes.
Lorsque le service de scoring n'est pas disponible, le système attend jusqu'à deux secondes avant d'appliquer un score par défaut. Cette limite évite que l'indisponibilité du service ne retarde excessivement la création d'un prospect.

Les mesures ont été effectuées sur un environnement de test simple. Elles montrent que la création d'un prospect reste possible même en l'absence du service de scoring.

Une amélioration possible consisterait à effectuer le calcul du score après la création du prospect, en arrière-plan, afin d'éviter toute attente supplémentaire en cas d'indisponibilité du service.

\subsection{Points forts}

\begin{description}[itemsep=1pt,parsep=0pt]
\item[Modules complémentaires.] La plateforme réunit dans un même environnement les fonctionnalités liées à la relation client et aux activités d'investissement.

\item[Architecture par microservices.] L'organisation retenue facilite l'ajout de nouvelles fonctionnalités et limite les dépendances entre les différents domaines métier.

\item[Gestion des accès et suivi des opérations.] Les mécanismes d'authentification, d'autorisation et de journalisation contribuent à sécuriser l'utilisation du système.

\item[Validation en environnement réel.] La plateforme a été déployée sur un serveur VPS AlmaLinux \cite{almalinux_official} afin de vérifier son fonctionnement et les procédures de mise à jour des services.
\end{description}



\subsection{Points faibles et limites}

\begin{description}[itemsep=1pt,parsep=0pt]
    \item[Périmètre encore incomplet.] La plateforme couvre un périmètre large,
          mais toutes les fonctionnalités et toutes les connexions avec les
          autres modules bancaires ne sont pas encore finalisées.
    \item[Données non officielles.] Les cotations BRVM \cite{brvm_actions}
          utilisées dans l'environnement de validation proviennent d'une source
          publique tierce \cite{brvm_data_public} mise à jour quotidiennement,
          et non directement du flux de marché officiel, dont l'intégration
          reste à prévoir pour une exploitation opérationnelle. De même,
          l'exécution des ordres est simulée, faute de connexion à un
          intermédiaire de marché réel.
    \item[Intégration de l'IA.] La pertinence métier des scores de prospect et
          des indicateurs de risque de départ relève du module IA et devra être
          confirmée sur des données réelles ; côté intégration, le parcours
          consomme déjà ces scores et gère l'indisponibilité du service.
    \item[Industrialisation du déploiement.] Le déploiement de validation repose sur Docker \cite{docker_docs} et Docker Compose \cite{docker_compose_docs} sur le VPS AlmaLinux \cite{almalinux_official}. Cette base reste à compléter par une supervision renforcée, des sauvegardes et une stratégie de reprise avant la production.
\end{description}

\subsection{Perspectives d'amélioration}

Les perspectives prioritaires concernent l'achèvement des intégrations avec les
autres modules, le raccordement à la BRVM \cite{brvm_actions} comme source officielle de marché, l'amélioration des
modèles de scoring et de churn, la stabilisation de la stratégie de déploiement
et l'ouverture de la plateforme vers des usages mobiles.

% ─────────────────────────────────────────────────

\section*{Conclusion}
\phantomsection
\addcontentsline{toc}{section}{Conclusion}

Ce chapitre a présenté les résultats obtenus avec I-SIB. La plateforme permet de centraliser les informations liées à la relation client, de suivre les prospects et de prendre en charge les principales fonctionnalités du module SGI.
Les tests menés au cours du projet ont permis de valider les principaux choix réalisés lors de la conception et du développement.

Les fonctionnalités mises en œuvre couvrent les besoins définis pour les modules étudiés. D'autres évolutions pourront être envisagées par la suite, notamment autour de l'exploitation de données réelles et de l'ajout de nouvelles fonctionnalités.

